\documentclass[12pt,a4paper]{article}
\usepackage[UTF8, heading=true]{ctex}
\usepackage{geometry}
\usepackage{multirow}
\usepackage{booktabs}
\usepackage{graphicx}
\usepackage{float}
% 页边距设置
\geometry{left=2.5cm,right=2.5cm,top=2.5cm,bottom=2.5cm}
% 字体与行距设置
\ctexset{
    section = { format=\bfseries\small },
    subsection = { format=\bfseries\small }
}
\linespread{1.5}
\fontsize{10.5pt}{12pt}\selectfont

\begin{document}
% 标题
\begin{center}
    \textbf{\small 智能垃圾分类处理系统仿真程序设计}
\end{center}

% 基本信息
\noindent
\textbf{成  绩}\hspace{4cm} \textbf{指导教师}：罗志祥 \\
\textbf{专业班级}：光电信息科学与工程202409 \hspace{1cm} \textbf{姓  名}：陈恪瑾 \hspace{1cm} \textbf{学  号}：U202413925 \\
\textbf{联系电话}：18271851291

\section{任务要求}
\subsection{实验目的}
\begin{itemize}
    \item 掌握使用 AT89S52 的 I/O 口；
    \item 理解 AT89S52 单片机的基本结构和工作原理；
    \item 掌握分支结构，条件转移程序设计方法；
    \item 具备调试程序的基本技能。
\end{itemize}

\subsection{实验内容}
在 Keil 环境中设计一个基于 AT89S52 单片机的智能垃圾分类处理系统，用户通过 P2 口的输入和 P1 口输出（上电复位初始状态为全 1），模拟小区垃圾回收终端系统的运行状态，系统需要实现以下功能：
\begin{itemize}
    \item 小区居民通过键入相应按键选择垃圾类型（分别用 P2.4~P2.7 端口模拟按键输入，低电平有效，P2.7=0、P2.6=0、P2.5=0、P2.4=0 分别代表“塑料纸张”、“玻璃金属”、“厨余垃圾”和“危害垃圾”四类）。
    \item P1.0~P1.3 四位代表垃圾桶盖开关控制，控制名为“塑料纸张”、“玻璃金属”、“厨余垃圾”和“危化品垃圾”四个垃圾桶盖，高电平为关闭状态，低电平为打开状态。
    \item 垃圾桶盖每次打开的时间缺省值为 8s，然后自动关闭，请采用循环程序配合 NOP 指令完成延时子程序，精度控制在 0.1s 以内。
\end{itemize}
请编写程序，实现用户键入，控制相应的垃圾桶盖打开和关闭。

\subsection{提高部分（申优选做）}
\begin{itemize}
    \item 系统初始化后，随机生成 4 个 0~9 的整数，置于寄存器 R0~R3 中表示四个垃圾桶的剩余容量。用户每次投递之后，容量相应减少。当投放垃圾超出剩余容量后，对应的标志位发生变化（自行定义），便于提醒物业及时清理，此时再次按键投放时，该桶盖不再打开避免溢出。
    \item 其他个性化设计。
\end{itemize}

\section*{二、设计思路}
本系统基于AT89S52单片机核心架构，采用模块化程序设计思想，围绕垃圾分类投放的核心需求，拆分功能模块实现全流程控制，整体设计逻辑如下：
\begin{enumerate}
    \item \textbf{系统初始化模块}：上电后完成核心状态配置，包括P1口置全1（所有桶盖默认关闭）、P2口置全1（配置为输入模式）；调用伪随机数生成子程序，为4个垃圾桶生成0-9的初始剩余容量；清零所有溢出告警标志位，完成系统初始状态复位。
    \item \textbf{按键检测与分支跳转模块}：主循环中持续读取P2口状态，通过取反、屏蔽低4位操作判断是否有按键按下；无按键时持续循环检测，有按键时通过位判断指令，跳转到对应垃圾类型的处理分支，实现多分支条件转移控制。
    \item \textbf{垃圾投放与容量管理模块}：按键触发后，先校验对应垃圾桶的剩余容量：若容量为0，直接置位对应溢出标志位，返回主循环，禁止开盖；若容量非0，拉低对应P1口引脚打开桶盖，调用8s延时子程序后拉高引脚关闭桶盖，同时将对应容量寄存器递减1，完成一次完整投放流程。
    \item \textbf{精准延时子程序}：基于12MHz晶振（1个机器周期=1$\mu$s），采用三层嵌套循环配合NOP指令，精准计算每个指令的机器周期，实现8s精准延时，精度控制在0.1s以内，满足桶盖自动关闭的时序要求。
    \item \textbf{伪随机数生成子程序}：采用互质步长累加取余算法，以片内RAM 30H单元为随机数种子，选用与10互质的7作为固定步长，通过累加后对10取余的方式，生成4个0-9范围内的伪随机数，实现垃圾桶初始容量的随机初始化。
\end{enumerate}

\section*{三、资源分配}
\subsection*{（一）I/O口硬件资源分配}
\begin{table}[H]
    \centering
    \small
    \begin{tabular}{|c|c|c|c|}
        \hline
        I/O引脚 & 方向 & 功能定义 & 有效电平 \\
        \hline
        P1.0 & 输出 & 塑料纸张垃圾桶盖控制 & 低电平打开，高电平关闭 \\
        \hline
        P1.1 & 输出 & 玻璃金属垃圾桶盖控制 & 低电平打开，高电平关闭 \\
        \hline
        P1.2 & 输出 & 厨余垃圾垃圾桶盖控制 & 低电平打开，高电平关闭 \\
        \hline
        P1.3 & 输出 & 危化品垃圾垃圾桶盖控制 & 低电平打开，高电平关闭 \\
        \hline
        P2.4 & 输入 & 危害垃圾投放按键 & 低电平有效 \\
        \hline
        P2.5 & 输入 & 厨余垃圾投放按键 & 低电平有效 \\
        \hline
        P2.6 & 输入 & 玻璃金属投放按键 & 低电平有效 \\
        \hline
        P2.7 & 输入 & 塑料纸张投放按键 & 低电平有效 \\
        \hline
    \end{tabular}
\end{table}

\subsection*{（二）寄存器资源分配}
\begin{table}[H]
    \centering
    \small
    \begin{tabular}{|c|c|}
        \hline
        寄存器 & 功能定义 \\
        \hline
        R0 & 存储塑料纸张垃圾桶剩余容量，初始值为0-9随机数 \\
        \hline
        R1 & 存储玻璃金属垃圾桶剩余容量，初始值为0-9随机数 \\
        \hline
        R2 & 存储厨余垃圾垃圾桶剩余容量，初始值为0-9随机数 \\
        \hline
        R3 & 存储危化品垃圾垃圾桶剩余容量，初始值为0-9随机数 \\
        \hline
        R5-R7 & 8s延时子程序三层循环计数寄存器，用于精准控制延时时长 \\
        \hline
        累加器A & 临时存储I/O口状态、随机数运算结果、数据中转 \\
        \hline
    \end{tabular}
\end{table}

\subsection*{（三）片内RAM资源分配}
\begin{table}[H]
    \centering
    \small
    \begin{tabular}{|c|c|}
        \hline
        存储单元 & 功能定义 \\
        \hline
        30H单元 & 伪随机数生成子程序的种子存储单元，用于连续生成0-9随机数 \\
        \hline
        20H.0 & 塑料纸张垃圾桶溢出告警标志位，置1表示容量用尽 \\
        \hline
        20H.1 & 玻璃金属垃圾桶溢出告警标志位，置1表示容量用尽 \\
        \hline
        20H.2 & 厨余垃圾垃圾桶溢出告警标志位，置1表示容量用尽 \\
        \hline
        20H.3 & 危化品垃圾垃圾桶溢出告警标志位，置1表示容量用尽 \\
        \hline
    \end{tabular}
\end{table}

\section*{四、流程图}
（请用Visio/Draw.io绘制后，插入\verb|\includegraphics|命令）

\section*{五、源代码}
\begin{verbatim}
; ========================================================
; 文件名称：SmartTrash.asm
; 处理器：AT89S52 (12MHz晶振)
; 实验内容：智能垃圾分类处理系统仿真程序
; ========================================================
        ORG 0000H
        LJMP MAIN           ; 上电跳转至主程序，避开中断向量区

        ORG 0030H           ; 主程序起始地址，避开中断入口
; ========================================================
; 主程序：系统初始化+主循环按键检测
; ========================================================
MAIN:
        ; 1. 系统I/O口初始化
        MOV P1, #0FFH       ; P1.0~P1.3 上电为全1，所有垃圾桶盖默认关闭
        MOV P2, #0FFH       ; P2口置全1，配置为输入模式，准备检测按键

        ; 2. 生成4个0~9伪随机数，作为四个垃圾桶初始剩余容量
        MOV 30H, #03H       ; 给30H单元设置初始种子，用于伪随机数生成
        LCALL GET_RAND_0_9
        MOV R0, A           ; R0 = 塑料纸张垃圾桶初始容量（0-9随机数）
        LCALL GET_RAND_0_9
        MOV R1, A           ; R1 = 玻璃金属垃圾桶初始容量（0-9随机数）
        LCALL GET_RAND_0_9
        MOV R2, A           ; R2 = 厨余垃圾垃圾桶初始容量（0-9随机数）
        LCALL GET_RAND_0_9
        MOV R3, A           ; R3 = 危化品垃圾垃圾桶初始容量（0-9随机数）

        ; 3. 溢出告警标志位初始化，全部清零
        CLR 20H.0           ; 20H.0 = 塑料纸张桶溢出标志
        CLR 20H.1           ; 20H.1 = 玻璃金属桶溢出标志
        CLR 20H.2           ; 20H.2 = 厨余垃圾桶溢出标志
        CLR 20H.3           ; 20H.3 = 危化品桶溢出标志

; 主循环：持续检测按键输入
MAIN_LOOP:
        MOV A, P2           ; 读取P2口完整状态
        CPL A               ; 按键低电平有效，取反后按下的按键对应位变为1
        ANL A, #0F0H        ; 屏蔽低4位，仅保留P2.4~P2.7按键相关位
        JZ MAIN_LOOP        ; 无按键按下（全0），返回循环继续检测

        ; 多分支查询，判断具体按下的按键，跳转对应处理分支
        JB ACC.7, THROW_PLASTIC   ; ACC.7对应P2.7，塑料纸张按键按下
        JB ACC.6, THROW_GLASS     ; ACC.6对应P2.6，玻璃金属按键按下
        JB ACC.5, THROW_KITCHEN   ; ACC.5对应P2.5，厨余垃圾按键按下
        JB ACC.4, THROW_HAZARD    ; ACC.4对应P2.4，危化品垃圾按键按下
        SJMP MAIN_LOOP             ; 无有效按键，返回主循环

; ========================================================
; 垃圾投放处理分支模块
; ========================================================
; 塑料纸张垃圾投放处理
THROW_PLASTIC:
        CJNE R0, #00H, OPEN_PLAS  ; 剩余容量不为0，跳转开盖流程
        SETB 20H.0                ; 容量为0，置位溢出标志，提醒物业清理
        SJMP MAIN_LOOP            ; 禁止开盖，返回主循环
OPEN_PLAS:
        CLR P1.0                  ; 拉低P1.0，打开塑料纸张桶盖
        LCALL DELAY_8S            ; 调用8s延时子程序，保持桶盖打开
        SETB P1.0                 ; 拉高P1.0，关闭桶盖
        DEC R0                    ; 投递完成，剩余容量减1
        SJMP MAIN_LOOP            ; 返回主循环，等待下一次按键

; 玻璃金属垃圾投放处理
THROW_GLASS:
        CJNE R1, #00H, OPEN_GLAS  ; 剩余容量不为0，跳转开盖流程
        SETB 20H.1                ; 容量为0，置位溢出标志
        SJMP MAIN_LOOP            ; 禁止开盖，返回主循环
OPEN_GLAS:
        CLR P1.1                  ; 拉低P1.1，打开玻璃金属桶盖
        LCALL DELAY_8S            ; 8s延时
        SETB P1.1                 ; 关闭桶盖
        DEC R1                    ; 剩余容量减1
        SJMP MAIN_LOOP

; 厨余垃圾投放处理
THROW_KITCHEN:
        CJNE R2, #00H, OPEN_KITC  ; 剩余容量不为0，跳转开盖流程
        SETB 20H.2                ; 容量为0，置位溢出标志
        SJMP MAIN_LOOP            ; 禁止开盖，返回主循环
OPEN_KITC:
        CLR P1.2                  ; 拉低P1.2，打开厨余垃圾桶盖
        LCALL DELAY_8S            ; 8s延时
        SETB P1.2                 ; 关闭桶盖
        DEC R2                    ; 剩余容量减1
        SJMP MAIN_LOOP

; 危化品垃圾投放处理
THROW_HAZARD:
        CJNE R3, #00H, OPEN_HAZA  ; 剩余容量不为0，跳转开盖流程
        SETB 20H.3                ; 容量为0，置位溢出标志
        SJMP MAIN_LOOP            ; 禁止开盖，返回主循环
OPEN_HAZA:
        CLR P1.3                  ; 拉低P1.3，打开危化品桶盖
        LCALL DELAY_8S            ; 8s延时
        SETB P1.3                 ; 关闭桶盖
        DEC R3                    ; 剩余容量减1
        SJMP MAIN_LOOP

; ========================================================
; 子程序名：DELAY_8S
; 功  能：实现8s精准延时，精度控制在0.1s以内
; 参  数：无
; 原  理：12MHz晶振下，1个机器周期=1us
;         内循环：250次 × 4us = 1ms
;         中层循环：200次 × 1ms = 200ms
;         外层循环：40次 × 200ms = 8000ms = 8s
; ========================================================
DELAY_8S:
        MOV R5, #40         ; 最外层循环，40次 × 200ms = 8s
D_OUTER:
        MOV R6, #200        ; 中层循环，200次 × 1ms = 200ms
D_MIDDLE:
        MOV R7, #250        ; 内层循环，250次 × 4us = 1ms
D_INNER:
        NOP                 ; 1个机器周期，1us
        NOP                 ; 1个机器周期，1us
        DJNZ R7, D_INNER    ; 2个机器周期，2us；单次内循环共4us
        DJNZ R6, D_MIDDLE   ; 中层循环跳转
        DJNZ R5, D_OUTER    ; 外层循环跳转
        RET                 ; 延时完成，子程序返回

; ========================================================
; 子程序名：GET_RAND_0_9
; 功  能：生成0-9范围内的伪随机数
; 参  数：无，返回值存于累加器A
; 原  理：选用与10互质的7作为步长，累加后对10取余，
;         保证生成0-9完整循环的伪随机数，无重复死区
; ========================================================
GET_RAND_0_9:
        MOV A, 30H          ; 读取当前随机数种子
        ADD A, #07H         ; 累加固定步长7
MOD_10: 
        CLR C               ; 清除借位标志
        SUBB A, #0AH        ; 减去10，实现取余运算
        JNC MOD_10          ; 无借位（A>=10），继续减10
        ADD A, #0AH         ; 有借位，补回多减的10，得到0-9余数
        MOV 30H, A          ; 更新随机数种子，用于下次生成
        RET                 ; 随机数存于A，子程序返回

        END                 ; 程序结束
\end{verbatim}

\section*{六、程序测试方法与结果}
采用keil软件进行仿真测试。

效果如下：

将P2.6调为低电平时，观察到P1.1也变为低电平，表示打开玻璃金属垃圾桶。

（此处插入对应图片：\verb|\includegraphics[width=10cm]{图片路径}|）

下图为延时，由光标可以看到大致为8s。

（此处插入对应图片：\verb|\includegraphics[width=10cm]{图片路径}|）

下图可以看到，当R0为0时（表示第一个垃圾桶已装满），此时调节P2.7=0，并不会使P1.0打开。

（此处插入对应图片：\verb|\includegraphics[width=10cm]{图片路径}|）

\section*{七、实验问题分析与对策}
（根据实验实际情况填写）

\vspace{2cm}
\noindent
本人承诺:本报告内容真实，无伪造数据，无抄袭他人成果。本人完全了解学校相关规定，如若违反，愿意承担其后果。

\vspace{1cm}
\noindent
签字：\hspace{3cm} 陈恪瑾 \hspace{3cm} 2026年4月26日

\end{document}